% 4.符号说明.tex —— 四、符号说明
\section{符号说明}
主要符号见表 \ref{tab:symbol}。物理时间在方程中以秒计，图表使用小时或分钟时另标单位；算法中的局部迭代量在出现处定义。

\begin{table}[H]
  \centering
  \caption{四个问题共用的主要符号及适用范围}
  \label{tab:symbol}
  \resizebox{\textwidth}{!}{%
  \begin{tabular}{llll}
    \toprule
    符号 & 含义 & 单位 & 取值或适用范围 \\
    \midrule
    $r,R,L$ & 径向位置、半径、长度 & m & $0\le r\le R$，$R=0.02$，$L=0.25$ \\
    $R(t)$ & 收缩半径（移动边界） & m & 问题四：附件 2，由 $0.02$ 收缩至 $\approx0.012$ \\
    $E$ & 归一化坐标 $r/R(t)$ & 无量纲 & 问题四 Landau 变换后的固定域 $[0,1]$ \\
    $t$ & 时间 & s & 问题一至 $1800$ s，问题二至 $72$ h \\
    $T(r,t)$ & 药材温度 & $^\circ$C & 待求 \\
    $T_K$ & 绝对温度 & K & $T_K=T+273.15$ \\
    $C(r,t)$ & 药材干基含水率 & kg/kg & 待求 \\
    $T_\infty(t),C_\infty(t)$ & 环境温度、水分边界参照值 & $^\circ$C；kg/kg & 附件 1，后两问在 $4$ h 后保持末值 \\
    $\rho$ & 药材密度 & kg/m$^3$ & 问题一 $820$；问题二、三 $\rho(C)$（附录 3）；问题四 $\rho(C)$（附录 4） \\
    $c_p$ & 比热容 & J/(kg$\cdot$K) & 问题一 $2600$；问题二、三 $c_p(C)$（附录 3）；问题四 $c_p(C)$（附录 4） \\
    $k$ & 导热系数 & W/(m$\cdot$K) & 问题一 $0.36$；问题二、三 $k(C)$（附录 3）；问题四 $k(C)$（附录 4） \\
    $h$ & 对流换热系数 & W/(m$^2\cdot$K) & $25$ \\
    $h_m$ & 对流传质系数 & m/s & $8\times10^{-7}$ \\
    $\alpha$ & 热扩散系数 $k/(\rho c_p)$ & m$^2$/s & 问题一 $1.689\times10^{-7}$ \\
    $D$ & 有效水分扩散系数 & m$^2$/s & 问题一 $D(C)$；问题二、三 $D(C,T)$（附录 3）；问题四 $D(C,T)$（附录 4） \\
    $\Delta r,N$ & 径向步长、区间数 & m；无量纲 & 均匀网格 $\Delta r=R/N$ \\
    $\Delta t,\Delta t_{\rm out}$ & 积分步长、输出时间间隔 & s & 问题三输出间隔 $60$ s \\
    $H,m$ & 宏观步长、宏观步内子步数 & s；无量纲 & 问题三时间推进控制量 \\
    $\theta$ & 时间加权参数 & 无量纲 & CN 取 $\theta=1/2$ \\
    $Fo$ & 网格 Fourier 数 $\alpha\Delta t/\Delta r^2$ & 无量纲 & 线性显式稳定性分析 \\
    $q$ & 分级网格参数 & 无量纲 & $q=1$ 对应均匀网格 \\
    $w_j$ & 径向控制体积积分权重 & m$^2$ & $\sum_jw_j=R^2/2$ \\
    $M$ & 径向加权含水量 $\sum_jw_jC_j$ & m$^2\cdot$kg/kg & 非直接物理总质量 \\
    $t_{\rm dry},t_{\rm rep}$ & 数值达标时刻、分钟报告时刻 & s & 以全场 $C<0.15$ 判定 \\
    \bottomrule
  \end{tabular}}
\end{table}

常热物性数值仅适用于问题一，问题二、三的状态函数见式 \eqref{eq:p2_param}、\eqref{eq:p2_D}，问题四的状态函数见式 \eqref{eq:p4_param}、\eqref{eq:p4_D}。本文以 $h$ 表示换热系数、$H$ 表示宏观步长；后文局部使用的步长整数倍数 $k$ 与导热系数 $k(C)$ 按上下文区分。
